\begin{document}
\maketitle

\begin{abstract}
This note is the thin entry document for the TFPT 4.5 series. It does not attempt to prove the
full theory. Its purpose is to state what TFPT claims, what it does not claim, how the closed branch
is organized, and where each load-bearing argument is isolated in the paper sequence.
\end{abstract}

\begin{tfptfrontbox}
\textbf{Inputs from previous papers.} None. This is the public orientation layer for the series.

\textbf{New theorem contribution.} None. The note contributes only organization, status discipline,
and a dependency map.

\textbf{Not claimed here.} No carrier proof, no exact electromagnetic calculation, no CMB fit, no
mass ledger, no E8 stage atlas, and no nonperturbative QFT proof.

\textbf{Falsification or audit surface.} The note can fail if it misstates the dependency order,
overstates the status of a downstream module, or hides where an assumption first enters.
\end{tfptfrontbox}


\begin{tfptclaimcontract}
\textbf{Claim.} Orientation and dependency map only; no theorem proof.\par
\textbf{Inputs.} None.\par
\textbf{First assumptions.} Status vocabulary and series dependency order.\par
\textbf{Proof status.} Reader contract / orientation.\par
\textbf{Kill condition.} Misstates dependency order or promotes downstream targets to core claims.\par
\end{tfptclaimcontract}

\begin{tfptstatusbox}
The TFPT 4.5 series uses a uniform \textbf{five-tier} status
taxonomy. Every load-bearing claim in every paper is tagged
inline using \stag{$\cdot$} with one of the six labels below.

\begin{description}[leftmargin=2em, itemsep=2pt]
    \item[\stag{L} \,Lean verified.]\ A theorem in the companion
        Lean 4 project
        \texttt{experiments/lean4-carrier-rigidity},
        signature-locked in \texttt{AuditContract.lean}, axioms
        only in
        $\{\texttt{propext},\ \texttt{Classical.choice},\ \texttt{Quot.sound}\}$.
    \item[\stag{I} \,Lean interfaced.]\ The Lean-side target
        type for the statement exists (a structure or theorem
        signature in \texttt{TfptCarrier}). The TFPT-internal
        proof producing an inhabitant of that type remains open.
    \item[\stag{M} \,Manuscript theorem.]\ A proof exists in
        the present manuscript or its source manuscript at
        standard mathematical-writing rigour. Not formalised
        in Lean.
    \item[\stag{S} \,Standard-theorem interface.]\ A routine
        consequence of a standard result in the literature
        (e.g.\ algebraic geometry of $\mathrm{Proj}$); cited
        rather than reproved.
    \item[\stag{C} \,Conditional downstream proposal.]\
        Mathematically formulated, but neither Lean verified,
        nor Lean interfaced, nor a present manuscript theorem.
        Conditional on upstream
        \stag{L}/\stag{I}/\stag{M}/\stag{S} claims.
    \item[\stag{P} \,Programmatic target.]\ A target for future
        work; not a claim. Stage-2 pipelines, comparison
        ledgers, sky-realisation tasks fall here.
\end{description}

\noindent
\textbf{Crucial distinction: \stag{L} \boldmath$\neq$\unboldmath\ \stag{M}.}
"Lean verified" and "proved in manuscript" are both valuable
but not identical. The label \stag{L} is reserved exclusively
for statements that elaborate inside the companion Lean project
and survive its signature-lock audit. Manuscript-level proofs
of comparable rigour are \stag{M}; do not promote them.

\noindent
Numerics depending on any \stag{C}/\stag{P} step are
\emph{conditional readouts}, not predictions.
\end{tfptstatusbox}

\tableofcontents

\section{What TFPT Claims}

TFPT is organized as a boundary-polarized spectral theory whose primitive input is a one-sided
boundary datum
\[
\Tbound=(\mathcal A_+,\mathcal H_+,D_+,J,\Gamma,B_\Sigma).
\]
The main theorem chain in the source draft is not a collection of unrelated numerological
readouts. It is a staged reconstruction:
\[
\Tbound
\Rightarrow
(\taudbl,\iC,\Pprim,[u_\Sigma],\cthree)
\Rightarrow
d^\star_{\mathrm{disc}}
\Rightarrow
\Padm
\Rightarrow
\mathcal M_{d^\star_{\mathrm{disc}}}
\Rightarrow
\mathcal Z_{\mathrm{cl}}=\{x^\star_{\mathrm{cl}}\}
\Rightarrow
\Tstar .
\]
The split series follows this burden of proof rather than the order in which a reader might find
the numerics exciting.

\section{What TFPT Does Not Claim at This Level}

The orientation note does not certify every downstream comparison row as a theorem. In particular,
the CMB interface, sky-map realization, horizon continuations, transient channels, and E8 scale
grammar remain outside the primitive theorem chain. They are comparison and programmatic layers
unless explicitly promoted by a later proof obligation.

\section{The Three Decoders}

The source manuscript condenses the closed branch through three decoders:
\[
\begin{aligned}
Y&\quad\text{generates structure},\\
[u_\Sigma]=1&\quad\text{generates counting},\\
u:=\phiz&\quad\text{generates bridge observables}.
\end{aligned}
\]
This orientation is the reason the paper series separates the primitive kernel, the carrier
packet, the precision readouts, the QFT closure layer, the metrology layer, and the cosmology
interface.

\section{Status Matrix}

\begin{center}
\footnotesize
\begin{tabularx}{\textwidth}{@{}p{0.24\textwidth}p{0.08\textwidth}p{0.22\textwidth}Y@{}}
\toprule
\textbf{Layer} & \textbf{Tier} & \textbf{Lean-side target / reference} & \textbf{Notes / location} \\
\midrule
Boundary primitive kernel: algebraic shadow
  & \stag{L}
  & \texttt{BoundaryPolarization}, \texttt{BoundaryPolarization.ofCarrierInvolution}
  & Algebraic equivalence between the polarisation picture and the involution picture is Lean-proved. Paper 1 + Lean note. \\
Boundary primitive kernel: analytic Paper-1 content
  & \stag{I}
  & target type \texttt{BoundaryPolarization}
  & Analytic Calder\'on derivation of $H = H_+ \oplus H_-$ is the upstream TFPT obligation. Paper 1 body \stag{M}/\stag{C}. \\
$3+2$ carrier algebra:
  $(q_-,q_+)=(-2,3),\ \tr Y=0,\ 6Y^2{-}Y{-}\mathbf 1{=}\mathbf 0$
  & \stag{L}
  & \texttt{CarrierPremises.hypercharge\_carrier\_packet}
  & Lean theorem, signature-locked. Paper 2 \S\ref{sec:formal-status} \& \Cref{thm:carrier-rigidity}. \\
Higgs rank-2 discharge (algebraic shadow $K[X_0,X_1]_1$)
  & \stag{L}
  & \texttt{HiggsTopForm}, \texttt{HiggsSchemeCohomologyShadow}
  & Lean-proved up to the algebraic shadow. \\
Higgs rank-2: Proj-side scheme-cohomology identification
  $H^0(\mathrm{Proj}\,K[X,Y], \mathcal O(1)) \simeq K[X,Y]_1$
  & \stag{S}
  & ---
  & Standard algebraic-geometry theorem; cited, not reproved. Paper 5 \stag{M} sketch. \\
Yukawa rank-3 forward (Stage D $\Rightarrow \dim E_-=3$)
  & \stag{L}
  & \texttt{YukawaTrilinearForm.finrank\_E\_eq\_three}
  & \\
Yukawa Stage-D characterisation:
  $\exists\,\omega\,(\textrm{CI}{\wedge}\textrm{CS}) \iff \dim E_-=3$
  & \stag{L}
  & \texttt{YukawaStageDExistence.stage\_D\_iff\_finrank\_eq\_three}
  & Both directions Lean-verified under \texttt{Nontrivial}\,$E_-$. \\
Yukawa: TFPT-canonical $\omega$ from boundary kernel
  & \stag{I}
  & \texttt{BoundaryYukawaKernel}
  & Existence Lean-closed; canonicity from TFPT data is open. \\
Seam winding
  $[u_\Sigma]=1 \Rightarrow$ primitive oriented determinant carrier
  & \stag{I}
  & \texttt{SeamWindingData}
  & Typed target available; topological derivation open. \\
Standard-Model derived representations,
  $\mathbb Z_6$ quotient, family count, occupancy, abelian index
  & \stag{C}
  & ---
  & Conditional downstream readouts. Paper 2 \stag{C} outlook (\Cref{sec:gauge-outlook,sec:packet-outlook,sec:counting-outlook}). \\
Electromagnetic / flavor readouts
  ($\alpha,\delta_{\mathrm{ph}},\lambda_C,\theta_{13}$, CKM, PMNS)
  & \stag{C}
  & ---
  & Closed-branch prediction layer. Paper 3. \\
Admissibility selector, strong CP, QFT reconstruction
  & \stag{C}
  & ---
  & Conditional nonperturbative closure under stated analytic hypotheses. Paper 4. \\
Boundary-normalised metrology
  & \stag{M}/\stag{C}
  & ---
  & Manuscript theorems on the metrology functor; downstream observables \stag{C}. Paper 5. \\
Cosmology interface (axion, reheating, leptogenesis, Stage-1 spectra)
  & \stag{C}
  & ---
  & Paper 6. \\
CMB operational closure / Planck pipeline (Stage-2 sky realisation)
  & \stag{P}
  & ---
  & Programmatic target with its own production-status ledger. Paper 7. \\
\bottomrule
\end{tabularx}
\end{center}

\medskip
\noindent
\textbf{Tier-collision warning.}\ The \emph{source manuscript}
(\texttt{../tfpt-42.tex}, extracted into \texttt{source\_extracts/})
uses the legacy vocabulary "theorem surface", "canonical rigid
object $\Tstar$", "closed theorem surface", and frequent
\stag{A}-labels in section headings such as
"\stag{A} Absolute spectral Planck closure". \emph{These legacy
\stag{A}-labels are not the \stag{L} of the present taxonomy.}
In source-extract bodies, every legacy \stag{A} should be read
as \stag{M} (manuscript theorem), and the "canonical rigid
object" / "closed theorem surface" language is the source-side
target picture, not the Lean-verified surface. The present
taxonomy is binding for the Frontmatter, claim contracts,
status boxes, and any new sections added in TFPT 4.5; the
source-extract bodies are draft material being progressively
re-tagged.

\noindent
The companion Lean note
(\texttt{experiments/lean4-carrier-rigidity/note\_carrier\_rigidity\_lean4.tex})
contains the full formalisation, the seven-check audit gate
(\texttt{AUDIT:~PASS}), and the signature-locked statement
contracts for every Status A theorem and every Status B
target type.

\section{Series Map}

\begin{enumerate}[leftmargin=1.5em]
    \item \textbf{Boundary Polarization and the Primitive Kernel of TFPT.}
    Establishes $\Tbound\Rightarrow(\taudbl,\iC,\Pprim,[u_\Sigma],\cthree)$ and keeps all
    Standard-Model claims out.
    \item \textbf{Carrier Rigidity, Hypercharge, and the Standard Model Packet from Boundary
    Polarization.}
    Establishes the rigid $3+2$ carrier, hypercharge, $S^+=\Lambda^{\mathrm{even}}E$,
    $G_{\mathrm{phys}}$, $N_{\mathrm{fam}}=3$, $\Omega_{\mathrm{adm}}=48$, $N_\Phi=1$, and
    $\bone=41/10$.
    \item \textbf{Electromagnetic Fixed Point and Flavor Transport on the Rigid TFPT Branch.}
    Isolates $\alpha$, $\delta_{\mathrm{ph}}$, $\lambda_C$, $\beta_{\mathrm{rad}}$, $\Omega_b$,
    $\theta_{13}$, CKM, and PMNS.
    \item \textbf{Admissibility, Strong CP, and Nonperturbative QFT Closure on the TFPT Branch.}
    Separates selector from dynamics and states the positivity, OS, local-net, scattering, and
    exact-flow obligations.
    \item \textbf{Geometric Hodge Closure and Dimensionless Metrology from the TFPT Boundary
    Branch.}
    Treats $\lambda_\Sigma$, Planck normalization, electroweak matching, and boundary-normalized
    observables.
    \item \textbf{Cosmology Interfaces of the TFPT Closed Branch.}
    Treats seam transfer, determinant-line phase, axion sector, reheating, leptogenesis, CMB
    spectra, and sky-realization targets.
\end{enumerate}

\section{Publication Order}

The recommended public order is Paper 0, Paper 2, Paper 1, Paper 3, Paper 4, Paper 5, and finally
Paper 6. The mathematical dependency order remains Paper 1 before Paper 2; the publication order
uses Paper 2 as the clearest visible hook while Paper 0 explains the dependency discipline.

\section{Source Boundary}

\begin{tfptsourcebox}
Primary source: \sourceDraft.

The orientation note draws especially from the abstract, the closed-architecture box, Sections 1
and 2, the canonical rigid-object discussion, the structural falsification map, and the conclusion.
It intentionally does not import the extended empirical ledgers, E8 stage atlas, CMB figures, or
mass tables.
\end{tfptsourcebox}

\clearpage
\section{Main Technical Development}

This section contains the main technical development assigned to this paper by the TFPT 4.5 clean split. Cross-paper background is referenced through dependency and scope boxes; extended backend material is kept in the Technical Companion.

\input{source_extracts/00_orientation_source.tex}


\begin{tfptexports}
Exports: series dependency order; status vocabulary; publication strategy; global falsification categories.
\end{tfptexports}

\end{document}
